\subsection{A positive comparison without logarithmic loss}
\label{band:section}

The comparison below works directly at a Frostman exponent, without
assuming that the energy at that exponent is finite. It also permits
nonconsecutive levels. All constants may depend on the fixed source and
pin geometry, their separation, and the indicated exponents, but are
independent of the levels, individual cell masses, and truncation height.

\begin{lemma}[Frequency-band comparison]\label{band:lemma}
Let \(1<s<2\), and let \(\lambda\) be a probability supported in
\(B(0,C_0\rho)\), satisfying
\(\lambda(B(z,t))\leq Kt^s\) for all \(z,t\).
Write \(F(\theta)\in L^2(\R)\) for its projected density, defined for
almost every \(\theta\in S^1\). Let \(\kappa\) be a finite positive measure
on a parameter space, with measurable angles \(\theta,\phi\) and a
measurable real shift \(v\), such that both angular marginals are bounded
by \(A\sigma\), where \(\sigma\) is normalized arclength. If
\(d(\theta,\phi)\leq\omega\) and \(|v|\leq v_0\), then
\begin{equation}\label{band:lemma-bound}
\int\|F(\theta)(\cdot-v)-F(\phi)\|_2^2\,d\kappa
\leq C_{s,C_0}AK\bigl((\rho\omega)^{s-1}+v_0^{s-1}\bigr).
\end{equation}
Neither the angles nor the shifts need be independent.
\end{lemma}
\begin{proof}
Finite first energy, which follows from the Frostman bound and \(s>1\),
gives the projected densities and their angular second integrability.
For \(R>0\), the positive Gaussian Fourier identity and the Frostman
bound give
\begin{equation}\label{band:gaussian}
\begin{aligned}
E_R(\lambda)&:=\int_{\R^2}e^{-|\xi|^2/R^2}
                   |\widehat\lambda(\xi)|^2\,d\xi\\
&=cR^2\iint e^{-c'R^2|x-x'|^2}\,d\lambda(x)d\lambda(x')
\leq C_sKR^{2-s}.
\end{aligned}
\end{equation}
For example, the last inequality follows by decomposing the spatial
integral into \(|x-x'|\leq R^{-1}\) and the annuli
\(2^jR^{-1}<|x-x'|\leq2^{j+1}R^{-1}\).
The same Gaussian identity for the signed coordinate measures gives
\begin{equation}\label{band:coordinate}
\int e^{-|\xi|^2/R^2}|\widehat{x_k\lambda}(\xi)|^2\,d\xi
\leq C_{C_0}\rho^2 E_R(\lambda)
\leq C_{s,C_0}K\rho^2R^{2-s}.
\end{equation}
Indeed the spatial Gaussian kernel is nonnegative and
\(|x_kx'_k|\leq C_0^2\rho^2\). This argument uses a Gaussian frequency
majorant; it does not assume a comparison between sharp annular
Fourier energies of \(x_k\lambda\) and \(\lambda\).

Choose smooth even functions \(\psi_R\), \(R\in2^{\mathbb Z}\),
supported on \(R/2\leq|\tau|\leq2R\), with bounded overlap and
\(\sum_R\psi_R(\tau)^2=1\) for \(\tau\ne0\). Set
$$
U_R(\theta)=\mathcal F^{-1}_\tau
                 [\psi_R(\tau)\widehat\lambda(\tau\theta)].
$$
This is a smooth function of \(\theta\) with values in \(L^2(\R)\).
Polar coordinates, \eqref{band:gaussian}, and
\eqref{band:coordinate}, at a fixed multiple of \(R\), imply
\begin{equation}\label{band:two-estimates}
\int\|U_R(\theta)\|_2^2\,d\sigma(\theta)\leq C_sKR^{1-s},
\qquad
\int\|\partial_\theta U_R(\theta)\|_2^2\,d\sigma(\theta)
\leq C_{s,C_0}K\rho^2R^{3-s}.
\end{equation}
For the second estimate, the angular derivative of
\(\widehat\lambda(\tau\theta)\) is a linear combination, with bounded
coefficients, of \(\tau\widehat{x_k\lambda}(\tau\theta)\).

The fundamental theorem of calculus along the shorter angular arc and
Cauchy--Schwarz give, when \(d(\theta,\phi)\leq\omega\),
$$
\|U_R(\theta)-U_R(\phi)\|_2^2
\leq \omega\int_{d(z,\theta)\leq\omega}
                         \|\partial_zU_R(z)\|_2^2\,dz.
$$
For \(\omega\) exceeding the circle diameter one may replace it by that
diameter. Integrating with the first angular marginal, using its bound
by \(A\sigma\), and applying Fubini bounds this by
\(CA\omega^2\|\partial_\theta U_R\|_{L^2_\theta L^2}^2\).
The trivial squared-difference bound, using both marginals, is
\(CA\|U_R\|_{L^2_\theta L^2}^2\). Consequently
\begin{equation}\label{band:angular}
\int\|U_R(\theta)-U_R(\phi)\|_2^2\,d\kappa
\leq CAK\min\{1,(\rho\omega R)^2\}R^{1-s}.
\end{equation}
Similarly, Plancherel and
\(|e^{-2\pi i\tau v}-1|^2\leq C\min\{1,(v_0R)^2\}\)
on the band imply
\begin{equation}\label{band:translation}
\int\|U_R(\theta)(\cdot-v)-U_R(\theta)\|_2^2\,d\kappa
\leq CAK\min\{1,(v_0R)^2\}R^{1-s}.
\end{equation}
The shift may depend arbitrarily on the parameter: this bound holds
before its integration.

For every parameter outside the angular marginal null sets, Plancherel
and \(\sum_R\psi_R^2=1\) identify the squared norm of the full difference
with the sum of its band squared norms. Scalar frequency localization
is preserved by the parameter-dependent shifts. Thus Tonelli, the
squared triangle inequality, \eqref{band:angular}, and
\eqref{band:translation} reduce the proof to
$$
\sum_{R\in2^{\mathbb Z}}\min\{1,(\delta R)^2\}R^{1-s}
\leq C_s\delta^{s-1}\qquad(\delta>0).
$$
Split the sum at \(R=\delta^{-1}\). The lower-frequency exponent is
\(3-s>0\), and the upper-frequency exponent is \(1-s<0\), so both sums
are geometric. A zero value of \(\delta\) contributes zero. This proves
\eqref{band:lemma-bound}, with no logarithmic factor.
\end{proof}

\begin{proposition}[Weighted comparison at arbitrary levels]
\label{band:comparison}
Use the separated probabilities, radial exponent \(q>1\), quantity
\(B\), fixed centers \(x_Q\), and positive affine densities \(f_n\)
from Section~\ref{sec:coherenttree}. Assume that the source \(\mu\) is
\(s\)-Frostman for some \(1<s<2\). For its positive-mass dyadic squares
at level \(m\), define
$$
\mathcal N_m=\left(\sum_{Q\text{ at level }m}\sqrt{p_Q}\right)^2,
\qquad
Z_{n,m}=2^{-m-2n(s-1)}\mathcal N_m.
$$
For all sufficiently large integers \(n<m\) and all \(A\geq1\),
\begin{equation}\label{band:comparison-bound}
\|f_m-f_n\|_{L^1(\nu\times dt)}
\leq C(AZ_{n,m})^{1/2}+CB A^{1-q}.
\end{equation}
In particular \(\mathcal N_m\leq N_m\), where \(N_m\) is the number of
occupied squares. There is no restriction \(m\leq2n\).
\end{proposition}
\begin{proof}
Write \(r=2^{-n}\), \(h=2^{-m}\). Let \(Q\) be a level-\(m\) square
and \(P\) its level-\(n\) ancestor. Apply both affine maps to the same
conditional probability \(\mu_Q\), recentered at \(b=x_Q\). Its support
lies in \(B(0,Ch)\), and its Frostman constant is at most \(C_\mu/p_Q\).
Put \(\theta=\Theta_{x_P}(y)\), \(\phi=\Theta_b(y)\). Separation gives
\(d(\theta,\phi)\leq Cr\), and the difference between the constant
terms of the two affine maps is
$$
v(y)=\theta\cdot(b-y)-|b-y|,
\qquad |v(y)|=|b-y|\,|\theta\cdot\phi-1|\leq Cr^2.
$$
Retain pins where both selected angular densities are at most \(A\).
Their coupled angular law has both marginals bounded by \(A\sigma\);
the discarded pin mass is at most \(A^{1-q}(F_P+F_Q)\).
Lemma~\ref{band:lemma} therefore bounds the retained integrated squared
second norm of the conditional affine-density difference by
$$
\frac{CA}{p_Q}\bigl((hr)^{s-1}+r^{2s-2}\bigr)
\leq \frac{CA}{p_Q}r^{2s-2},
$$
since \(h\leq r\).

For each pin each of the two conditional densities is supported on an
interval of length at most \(Ch\). Their union has length at most
\(2Ch\), even if the gap between those intervals exceeds \(h\).
Cauchy--Schwarz on this union and then over retained pins gives the
conditional first-norm bound
\(CA^{1/2}h^{1/2}r^{s-1}p_Q^{-1/2}\).
On discarded pins the first norm is at most two. Multiply by \(p_Q\)
and sum all level-\(m\) squares. The retained term is
$$
CA^{1/2}h^{1/2}r^{s-1}\sum_Q\sqrt{p_Q}
=C(AZ_{n,m})^{1/2}.
$$
The discarded terms sum to at most \(CB A^{1-q}\), because
$$
\sum_Qp_QF_Q\leq2B,
\qquad
\sum_{Q\text{ at level }m}p_QF_{P(Q)}
=\sum_{P\text{ at level }n}p_PF_P\leq2B.
$$
Decomposing every parent source into its descendants now proves
\eqref{band:comparison-bound}. Finally
\((\sum_Q\sqrt{p_Q})^2\leq N_m\sum_Qp_Q=N_m\).
\end{proof}

\begin{theorem}[Selected-level positive convergence]
\label{band:criterion}
Under the hypotheses of Proposition~\ref{band:comparison}, let
\(n_j\) be strictly increasing integers tending to infinity, and put
\(\eta_q=(q-1)/(2q-1)\). If
\begin{equation}\label{band:criterion-series}
\sum_j\left[
2^{-n_{j+1}-2n_j(s-1)}\mathcal N_{n_{j+1}}
\right]^{\eta_q}<\infty,
\end{equation}
then the raw pinned law \((d_y)_*\mu\) is absolutely continuous for
\(\nu\)-almost every \(y\).
\end{theorem}
\begin{proof}
Enlarge \(B\) to be at least one. Write \(Z_j=Z_{n_j,n_{j+1}}\), and
omit finitely many terms so \(Z_j\leq1\). In
\eqref{band:comparison-bound} take
\(A_j=(B/\sqrt{Z_j})^{1/(q-1/2)}\). Both resulting terms are at most
\(CB^{1/(2q-1)}Z_j^{\eta_q}\). Thus \(f_{n_j}\) converges in
\(L^1(\nu\times dt)\) to a nonnegative density of mass one. The uniform
affine approximation error \(C2^{-2n_j}\) identifies the limit measure
as the pushforward of \(\mu\times\nu\) by
\((x,y)\mapsto(y,|x-y|)\). Uniqueness of disintegration proves the
assertion about the actual positive distance probabilities.
\end{proof}

For consecutive levels the count-only sufficient condition is
\begin{equation}\label{band:consecutive}
\sum_n\left[2^{-n(2s-1)}N_{n+1}\right]^{\eta_q}<\infty.
\end{equation}
In particular \(N_n\leq C2^{n(2s-1)}n^{-a}\), with
\(a>1/\eta_q\), suffices. This improves the preceding variable-exponent
endpoint estimate by removing its factor of \(n\). A dimension equality
by itself does not imply \eqref{band:consecutive}; selected levels can
instead be tested using \eqref{band:criterion-series}.

\subsection{The exact energy at the curvature scale}
\label{band:energy-section}

There is a stronger spatial criterion before the Frostman estimate is
inserted. For a finite positive measure \(\lambda\), define
\begin{equation}\label{band:truncated-energy}
J_\delta(\lambda)=\iint k_\delta(x-x')\,d\lambda(x)d\lambda(x'),
\qquad
k_\delta(z)=\min\{|z|^{-1},\delta^2|z|^{-3}\},
\end{equation}
with value \(+\infty\) on the diagonal. This is a positive spatial
kernel and is bounded by the first-energy kernel.

\begin{lemma}[Positive spatial form]\label{band:spatial}
For any finite positive measure of finite first energy and any fixed
\(c>0\),
\begin{equation}\label{band:kernel-equivalence}
\sum_{R\in2^{\mathbb Z}}
\min\{1,(\delta R)^2\}R^{-1}E_{cR}(\lambda)
\asymp_c J_\delta(\lambda).
\end{equation}
In Lemma~\ref{band:lemma}, the Frostman hypothesis can consequently be
replaced by finite first energy, with the right side replaced by
\(CA[J_{\rho\omega}(\lambda)+J_{v_0}(\lambda)]\). A term with zero
scale parameter is interpreted as zero.
\end{lemma}
\begin{proof}
The Gaussian identity and Tonelli express the sum in
\eqref{band:kernel-equivalence} as a positive double integral with kernel
$$
\sum_{R\in2^{\mathbb Z}} C_cR
                  \min\{1,(\delta R)^2\}e^{-c_cR^2|z|^2}.
$$
Write \(u=|z|>0\). If \(u\leq\delta\), its upper bound is
\(C\sum_R Re^{-c_cR^2u^2}\leq C/u\). If \(u>\delta\), use
\(\min\{1,(\delta R)^2\}\leq\delta^2R^2\), obtaining
\(C\delta^2\sum_R R^3e^{-c_cR^2u^2}\leq C\delta^2/u^3\).
For the matching lower bound, retain a dyadic \(R\) comparable to
\(u^{-1}\). The diagonal has zero product mass under finite first
energy. These observations prove \eqref{band:kernel-equivalence}.

Before using the Frostman upper bound, the two estimates in
\eqref{band:two-estimates} are bounded respectively by
\(CR^{-1}E_{cR}(\lambda)\) and
\(C\rho^2R E_{cR}(\lambda)\). The coordinate estimate here uses only
the first inequality of \eqref{band:coordinate}. Repeating the
bandwise coupled-angle and scalar-translation proof thus gives the
two sums in \eqref{band:kernel-equivalence}, with parameters
\(\rho\omega\) and \(v_0\). This proves the refinement.
\end{proof}

\begin{proposition}[Curvature-scale energy criterion]
\label{band:energy-criterion}
Suppose that the separated compact source and pin probabilities have
the fixed radial data \(q>1,B<\infty\) and centers used above, and that
\(I_1(\mu)<\infty\). Define
\begin{equation}\label{band:energy-coefficient}
W_{n,m}=2^{-m}\left[
\sum_{Q\text{ at level }m}p_Q
           J_{2^{-2n}}(\mu_Q)^{1/2}\right]^2,
\qquad n<m.
\end{equation}
Then
$$
\|f_m-f_n\|_1\leq C(AW_{n,m})^{1/2}+CB A^{1-q}
\qquad(A\geq1).
$$
For any increasing sequence \(n_j\), the condition
\(\sum_j W_{n_j,n_{j+1}}^{\eta_q}<\infty\) therefore implies absolute
continuity of the raw pinned law for \(\nu\)-almost every pin.
\end{proposition}
\begin{proof}
Use Lemma~\ref{band:spatial} in the proof of
Proposition~\ref{band:comparison}. The angular parameter is at most
\(Chr\), and the translation parameter is at most \(Cr^2\), with
\(h=2^{-m}\leq r=2^{-n}\). Monotonicity of \(J_\delta\) in
\(\delta\), and \(J_{C\delta}\leq\max(1,C^2)J_\delta\), give the
conditional squared second-norm bound \(CAJ_{r^2}(\mu_Q)\).
Conversion on the union of the two length-\(Ch\) supports gives
\(C\sqrt{Ah}\,J_{r^2}(\mu_Q)^{1/2}\). The same weighted sum and bad-pin
bound prove the claimed first-norm estimate. Optimization in \(A\),
positive first-norm convergence, and identification of the limiting
law are exactly as in Theorem~\ref{band:criterion}.
\end{proof}

For an \(s\)-Frostman source, a direct decomposition of the kernel at
\(|x-x'|=\delta\) gives
$$
J_\delta(\mu_Q)\leq \frac{C_sC_\mu}{p_Q}\delta^{s-1},
$$
recovering \(W_{n,m}\leq C Z_{n,m}\). The exact criterion retains more
of the source geometry than this covering bound. It still retains all
close source pairs: for every conditional probability,
$$
J_\delta(\mu_Q)\geq
\delta^{-1}(\mu_Q\times\mu_Q)
             \{(x,x'):|x-x'|\leq\delta\}.
$$
Thus replacing the energy by its curvature-scale form does not
automatically make concentrated near-source pairs harmless. Neither
criterion asserts that a new dimension-only region or a bare critical
dimension equality follows without a further geometric estimate.
